\subsection{模的加性集的格罗滕迪克群}

设 A 是一个环。在本小节中，我们考虑 A-模的类的一个加性集 C；我们将 C 与自由交换群 \(\mathbf{Z}^{(\mathcal{C})}\) 的典范基等同。对任何由 C 型模组成的正合列

\[
0 \longrightarrow \mathrm{E} ^ {\prime} \longrightarrow \mathrm{E} \longrightarrow \mathrm{E} ^ {\prime \prime} \longrightarrow 0\tag{\( \mathcal{E} \)}
\]

我们用 \(r_{\mathcal{E}}\) 表示元素 \(\mathrm{cl}(\mathrm{E}) - \mathrm{cl}(\mathrm{E}') - \mathrm{cl}(\mathrm{E}'')\)（\(\mathbf{Z}^{(\mathcal{C})}\) 中的）。设 R 是 \(\mathbf{Z}^{(\mathcal{C})}\) 中由形如 \(r_{\mathcal{E}}\) 的元素生成的子群；商群 \(\mathbf{Z}^{(\mathcal{C})}/\mathrm{R}\) 称为 C 的格罗滕迪克群，记作 \(\mathrm{K}(\mathcal{C})\)。对一个 C 型模 E，我们将 \(\mathrm{cl}(\mathrm{E})\) 在 \(\mathrm{K}(\mathcal{C})\) 中的象记作 \([\mathrm{E}]_{\mathcal{C}}\)（当关于 C 可能有歧义时，有时也写作 {[}E{]}）。我们则有如下的泛性质。

命题 4. —— a) 映射 \(E \mapsto [E]_{\mathcal{C}}\)（从 C 到 \(\mathrm{K}(\mathcal{C})\)）是加性的。

\begin{enumerate}
\def\labelenumi{\alph{enumi})}
\setcounter{enumi}{1}
\tightlist
\item
  设 G 是一个交换群，\(\varphi : \mathcal{C} \to G\) 是模的一个加性函数。存在唯一的同态 \(u: K(\mathcal{C}) \to G\)，使得 \(\varphi(E) = u([E]_{\mathcal{C})}\) 对每个 \(\mathcal{C}\) 型模 E 成立。
\end{enumerate}

断言 a) 是显然的。我们来证明 b)。存在一个同态 \(u^{\prime}: \mathbf{Z}^{(\mathcal{C})} \to \mathrm{G}\)，它延拓 \(\varphi\)。由于 \(\varphi\) 是加性的，我们有 \(u^{\prime}(r_{\mathcal{E}}) = 0\)，对每个由 C 型模组成的正合列（ \(E\) ）皆然；因此 R 含于 \(u^{\prime}\) 的核。于是，在过渡到商群时，\(u^{\prime}\) 定义了从 \(\mathrm{K}(\mathcal{C})\) 到 G 的一个同态 u，并且显然对每个 C 型 A-模都有 \(\varphi(\mathrm{E}) = u([\mathrm{E}]_{\mathcal{C}})\)。群 \(\mathrm{K}(\mathcal{C})\) 由集合中的诸元素 \([E]_{C}\)（E 取遍 C）生成；u 的唯一性由此得出。

设 C 与 D 是 A-模的类的加性集，且 \(C \subset D\)。由于从 C 到格罗滕迪克群 K(D) 的映射 \(E \mapsto [E]_{D}\) 是加性的，存在一个同态 \(\gamma_{\mathcal{D},\mathcal{C}} : K(\mathcal{C}) \to K(\mathcal{D})\)，称为典范的，它由公式 \(\gamma_{\mathcal{D},\mathcal{C}}([E]_{\mathcal{C})} = [E]_{\mathcal{D}}\)（对每个 C 型模 E）刻画。它并不总是单的（VIII, 第 210 页，习题 13）。

例. —— *设 A 是一个诺特交换环，\(\Sigma\) 是它的谱。对任何整数 \(n \geqslant 0\)，用 \(\mathcal{C}^{\geqslant n}\) 表示其支集的余维数 \(\geqslant n\)（在 \(\Sigma\) 中）的有限生成 A-模的类所成之集。令 \(\mathrm{K}(n,\mathrm{A}) = \mathrm{K}(\mathcal{C}^{\geqslant n})\)，\(\gamma_{n} = \gamma_{\mathcal{C}^{\geqslant n},\mathcal{C}^{\geqslant n + 1}}\)。我们有一列同态

\[
\gamma_ {n}: \mathrm{K} (n + 1, \mathrm{A}) \longrightarrow \mathrm{K} (n, \mathrm{A}).
\]

我们可以证明（AC, VIII, §1, n° 5, 第 13 页，命题 10），在 K(n, A) 中，诸元素 \([A/p]_{C_{n}}\)（其中 p 取遍 \(\Sigma\) 中高度为 n 的元素所成之集）构成 \(\gamma_{n}\) 的象的一个补 Z-模的基。更精确地说，对每个 \(C \geqslant n\) 型模 E，我们有

\[
[ \mathrm{E} ] _ {\mathcal {C} ^ {\geqslant n}} \equiv \sum_ {\{\mathfrak {p} \in \Sigma | \operatorname{ht} (\mathfrak {p}) = n \}} \operatorname{long} _ {\mathrm{A} _ {\mathfrak {p}}} (\mathrm{E} _ {\mathfrak {p}}) \cdot [ \mathrm{A} / \mathfrak {p} ] _ {\mathcal {C} ^ {\geqslant n}} \quad (\mathrm{mod} \operatorname{Im} \gamma_ {n}). _ {*}
\]

群 \(\mathrm{K}(\mathcal{C})\) 由形如 \([E]_{\mathcal{C}}\)（\(E \in C\)）的元素生成，且由关系式 (1)（VIII, 第 184 页）我们有 \([E \oplus E']_{\mathcal{C}} = [E]_{\mathcal{C}} + [E']_{\mathcal{C}}\)。因此 \(\mathrm{K}(\mathcal{C})\) 的每个元素都形如 \([E]_{\mathcal{C}} - [F]_{\mathcal{C}}\)，其中 E 与 F 属于 C。

若 \(\mathrm{K}(\mathcal{C})\) 的一个元素对某个 C 型 A-模 E 形如 \([E]_{\mathcal{C}}\)，则称它是有效的。\(\mathrm{K}(\mathcal{C})\) 的有效元素所成之集记作 \(\mathrm{K}(\mathcal{C})^{+}\)；它是 \(\mathrm{K}(\mathcal{C})\) 的一个子幺半群，且 \(\mathrm{K}(\mathcal{C})\) 可与 \(\mathrm{K}(\mathcal{C})^{+}\) 的差群等同（I, §2, No.~4, 第 20 页）。

命题 5. —— 设 E 与 F 是 C 型模。等式 \([E]_{C} = [F]_{C}\) 成立当且仅当存在 C 型模的正合列

(ε)

\[
0 \longrightarrow \mathrm{L} \longrightarrow \mathrm{P} \longrightarrow \mathrm{M} \longrightarrow 0,\tag{\( \mathcal{F} \)}
\]

\[
0 \longrightarrow \mathrm{L} \longrightarrow \mathrm{Q} \longrightarrow \mathrm{M} \longrightarrow 0
\]

使得 \(\mathrm{E} \oplus \mathrm{Q}\) 同构于 \(\mathrm{F} \oplus \mathrm{P}\)。

所述条件是充分的，因为它蕴含关系式

\[
[ \mathrm{P} ] _ {\mathcal {C}} = [ \mathrm{L} ] _ {\mathcal {C}} + [ \mathrm{M} ] _ {\mathcal {C}}, \qquad [ \mathrm{Q} ] _ {\mathcal {C}} = [ \mathrm{L} ] _ {\mathcal {C}} + [ \mathrm{M} ] _ {\mathcal {C}}, \qquad [ \mathrm{E} ] _ {\mathcal {C}} + [ \mathrm{Q} ] _ {\mathcal {C}} = [ \mathrm{F} ] _ {\mathcal {C}} + [ \mathrm{P} ] _ {\mathcal {C}},
\]

由此得出 \([\mathrm{E}]_{\mathcal{C}} = [\mathrm{F}]_{\mathcal{C}}\)

现在假设我们有 \([E]_{C} = [F]_{C}\)。由群 \(\mathrm{K}(\mathcal{C})\) 的构造，存在 C 型模的正合列的两个有限族

\[
0 \longrightarrow \mathrm{G} _ {i} ^ {\prime} \longrightarrow \mathrm{G} _ {i} \longrightarrow \mathrm{G} _ {i} ^ {\prime \prime} \longrightarrow 0\tag{\((\mathcal{G}_i)\)}
\]

（\(i\in \mathrm{I}\)）与

\[
0 \longrightarrow \mathrm{H} _ {j} ^ {\prime} \longrightarrow \mathrm{H} _ {j} \longrightarrow \mathrm{H} _ {j} ^ {\prime \prime} \longrightarrow 0\tag{\((\mathcal{H}_j)\)}
\]

（\(j\in \mathrm{J}\)），使得在 \(\mathbf{Z}^{(\mathcal{C})}\) 中我们有

\[
\operatorname{cl} (\mathrm{E}) - \operatorname{cl} (\mathrm{F}) = \sum_ {j \in \mathrm{J}} r _ {\mathcal {H} _ {j}} - \sum_ {i \in \mathrm{I}} r _ {\mathcal {G} _ {i}}.
\]

更明确地，这个关系式可写成

\[
\begin{array}{c} \operatorname{cl} (\mathrm{E}) + \sum_ {i \in \mathrm{I}} \operatorname{cl} (\mathrm{G} _ {i}) + \sum_ {j \in \mathrm{J}} \operatorname{cl} (\mathrm{H} _ {j} ^ {\prime}) + \sum_ {j \in \mathrm{J}} \operatorname{cl} (\mathrm{H} _ {j} ^ {\prime \prime}) \\ = \operatorname{cl} (\mathrm{F}) + \sum_ {i \in \mathrm{I}} \operatorname{cl} (\mathrm{G} _ {i} ^ {\prime}) + \sum_ {i \in \mathrm{I}} \operatorname{cl} (\mathrm{G} _ {i} ^ {\prime \prime}) + \sum_ {j \in \mathrm{J}} \operatorname{cl} (\mathrm{H} _ {j}). \end{array}\tag{7}
\]

令 \(G = \bigoplus_{i \in I} G_i\)，\(G' = \bigoplus_{i \in I} G_i'\)，等等。通过过渡到直和，我们得到正合列

\begin{enumerate}
\def\labelenumi{(\Alph{enumi})}
\setcounter{enumi}{6}
\tightlist
\item
\end{enumerate}

\[
0 \longrightarrow \mathrm{G} ^ {\prime} \stackrel {p} {\longrightarrow} \mathrm{G} \stackrel {q} {\longrightarrow} \mathrm{G} ^ {\prime \prime} \longrightarrow 0,\tag{H}
\]

\[
0 \longrightarrow \mathrm{H} ^ {\prime} \stackrel {r} {\longrightarrow} \mathrm{H} \stackrel {s} {\longrightarrow} \mathrm{H} ^ {\prime \prime} \longrightarrow 0
\]

它们由 \(\mathcal{C}\) 型模组成。

其次，设 \(M_{1},\ldots,M_{m},N_{1},\ldots,N_{n}\) 是 C 型 A-模。若我们有 \(\sum_{i=1}^{m}\operatorname{cl}(M_{i})=\sum_{j=1}^{n}\operatorname{cl}(N_{j})\)（在群 \(\mathbf{Z}^{(\mathcal{C})}\) 中），则我们有 m=n，且存在一个置换 \(\sigma\in\mathfrak{S}_{m}\)，使得 \(\operatorname{cl}(M_{i})=\operatorname{cl}(N_{\sigma(i)})\) 对每个 \(1\leqslant i\leqslant m\) 成立（I, §7, No.~9, 第 95 页，命题 11）。因此，模 \(\bigoplus_{i=1}^{m}M_{i}\) 与 \(\bigoplus_{j=1}^{n}N_{j}\) 同构。特别地，由 (7) 我们推出存在一个从 \(E\oplus Q\) 到 \(F\oplus P\) 的同构，其中我们已令

\[
\mathrm{P} = \mathrm{G} ^ {\prime} \oplus \mathrm{G} ^ {\prime \prime} \oplus \mathrm{H}, \quad \mathrm{Q} = \mathrm{G} \oplus \mathrm{H} ^ {\prime} \oplus \mathrm{H} ^ {\prime \prime}.
\]

我们又令

\[
\mathrm{L} = \mathrm{G} ^ {\prime} \oplus \mathrm{H} ^ {\prime}, \quad \mathrm{M} = \mathrm{G} ^ {\prime \prime} \oplus \mathrm{H} ^ {\prime \prime}.
\]

模 L、M、P、Q 是 C 型的，且序列

\[
0 \longrightarrow \mathrm{L} \stackrel {\lambda} {\longrightarrow} \mathrm{P} \stackrel {\mu} {\longrightarrow} \mathrm{M} \longrightarrow 0,\tag{\( \mathcal{E} \)}
\]

（其中我们由下式定义 \(\lambda\) 与 \(\mu\)）

\[
\lambda (g ^ {\prime}, h ^ {\prime}) = \left(g ^ {\prime}, 0, r (h ^ {\prime})\right), \quad \mu (g ^ {\prime}, g ^ {\prime \prime}, h) = \left(g ^ {\prime \prime}, s (h)\right),
\]

是正合的。我们用同样的方式构造一个正合列

\[
0 \longrightarrow \mathrm{L} \longrightarrow \mathrm{Q} \longrightarrow \mathrm{M} \longrightarrow 0\tag{\( \mathcal{F} \)}
\]

；这就完成了证明。

集 C 对合成律 \((\mathrm{E}, \mathrm{E}^{\prime}) \mapsto \mathrm{cl}(\mathrm{E} \oplus \mathrm{E}^{\prime})\) 是一个交换幺半群。我们有时将交换幺半群 \(\mathcal{C}\)（I, §2, No.~4, 第 20 页）的差群记作 \(\mathrm{K}'(\mathcal{C})\)，并称之为 \(\mathcal{C}\) 的直和格罗滕迪克群。对任何 \(\mathcal{C}\) 型模 E，我们将 \(\operatorname{cl}(\mathrm{E})\) 在 \(\mathrm{K}'(\mathcal{C})\) 中的象记作 \([\mathrm{E}]_{\mathcal{C}}'\)。

命题 6. —— a) 映射 \(E \mapsto [E]_{\mathcal{C}}'\)（从 C 到 \(K'(\mathcal{C})\)）是模的一个弱加性函数。

\begin{enumerate}
\def\labelenumi{\alph{enumi})}
\setcounter{enumi}{1}
\item
  设 G 是一个交换群，\(\varphi : \mathcal{C} \to G\) 是模的一个弱加性函数。存在唯一的群同态 \(u: K'(\mathcal{C}) \to G\)，使得 \(\varphi(E) = u([E]_{\mathcal{C}}')\) 对每个 \(\mathcal{C}\) 型模 E 成立。
\item
  设 \(\mathrm{E}\) 与 \(\mathrm{F}\) 是 \(\mathcal{C}\) 型模。我们有 \([\mathrm{E}]_{\mathcal{C}}' = [\mathrm{F}]_{\mathcal{C}}'\) 当且仅当存在一个模 \(\mathrm{M}\)（\(\mathcal{C}\) 型的），使得 \(\mathrm{E} \oplus \mathrm{M}\) 同构于 \(\mathrm{F} \oplus \mathrm{M}\)。
\end{enumerate}

断言 a) 是显然的。断言 b) 由（I, §2, No.~4, 第 19 页，定理 1）得出，断言 c) 由（I, §2, No.~4, 第 18 页，命题 6）得出。

由于映射 \(E \mapsto [E]_{\mathcal{C}}\)（从 C 到 \(\mathrm{K}(\mathcal{C})\)）是模的一个弱加性函数，我们可以由它导出一个同态 \(u : \mathrm{K}'(\mathcal{C}) \to \mathrm{K}(\mathcal{C})\)。这个同态是满的，但并不总是同构（VIII, 第 191 页，注 2）。

设 \(R'\) 是 \(\mathbf{Z}^{(\mathcal{C})}\) 中由形如 \(r_{\mathcal{E}}\) 的元素（其中 \(\mathcal{E}\) 是 \(\mathcal{C}\) 型 A-模的一个分裂正合列）生成的子群，即由形如 \(\operatorname{cl}(\operatorname{E}' \oplus \operatorname{E}'') - \operatorname{cl}(\operatorname{E}') - \operatorname{cl}(\operatorname{E}'')\) 的元素（其中 \(E'\) 与 \(E''\) 是 \(\mathcal{C}\) 型模）生成的子群。从 \(\mathcal{C}\) 到商群 \(\mathbf{Z}^{(\mathcal{C})}/R'\) 的典范映射可延拓为群同态 \(v: K'(\mathcal{C}) \to \mathbf{Z}^{(\mathcal{C})}/R'\)。这是一个同构。事实上，从 \(\mathcal{C}\) 到 \(K'(\mathcal{C})\) 的典范映射可延拓为从 \(\mathbf{Z}^{(\mathcal{C})}\) 到 \(K'(\mathcal{C})\) 的一个群同态，其核包含 \(R'\)，因而通过过渡到商群可延拓为同态 \(v': \mathbf{Z}^{(\mathcal{C})}/R' \to K'(\mathcal{C})\)；显然 \(v\) 与 \(v'\) 是互逆的双射。

若 \(\mathrm{K}^{\prime}(\mathcal{C})\) 的一个元素对某个 C 型 A-模 E 形如 \([E]_{\mathcal{C}}^{\prime}\)，则称它是有效的。\(\mathrm{K}^{\prime}(\mathcal{C})\) 的有效元素所成之集记作 \(\mathrm{K}^{\prime}(\mathcal{C})^{+}\)。

